\documentclass{jsarticle}
\usepackage{ascmac,amsmath,amssymb,amsthm,bm,physics,arydshln,mathcomp}
\newtheorem*{answer*}{解答}

\begin{document}

\begin{center}
  数学科教育論1 \ \ \ \ 第三回課題
\end{center}

\begin{flushright}
  6122111 三森誠真
\end{flushright}

\begin{itembox}[l]{問3.8}
二項定理$(a+b)^n = \displaystyle \sum_{k=0}^n {}_n C_k \ a^k \ b^{n-k}$ が成り立つことを, 数学的帰納法を用いて証明せよ.
\end{itembox}

\begin{answer*}
\end{answer*}
(i) $n=1$ のとき, \ (左辺) = $a+b$.

\ \ \ \ \ \ \ \ \ \ (右辺) = ${}_1 C_0 \ b + {}_1 C_1 \ a = a+b$
なので成り立つ.

(ii) $n=k$ のとき, \ 式$(a+b)^k = \displaystyle \sum_{l=0}^k {}_k C_l \ a^l \ b^{k-l}$が成り立つと仮定すると, $n=k+1$ のとき,

\begin{align*}
(a+b)^{k+1} 
&= a(a+b)^k + b(a+b)^k \\
&= \displaystyle \sum_{l=0}^k {}_k C_l \ a^{l+1} \ b^{k-l} + \displaystyle \sum_{l=0}^k {}_k C_l \ a^l \ b^{k+1-l} \\
&= \displaystyle \sum_{s=1}^{k+1} {}_k C_{s-1} \ a^s \ b^{k+1-s} + \displaystyle \sum_{l=0}^k {}_k C_l \ a^l \ b^{k+1-l} \ (第一項にてl+1=s と置換)\\
&= \displaystyle \sum_{l=1}^{k+1} {}_k C_{l-1} \ a^l \ b^{k+1-l} + \displaystyle \sum_{l=0}^k {}_k C_l \ a^l \ b^{k+1-l} \\
&= a^{k+1} + \displaystyle \sum_{l=1}^{k} {}_k C_{l-1} \ a^l \ b^{k+1-l} + \displaystyle \sum_{l=1}^k {}_k C_l \ a^l \ b^{k+1-l} + b^{k+1} \\
&= \displaystyle \sum_{l=1}^{k} {}_{k+1} C_l \ a^l \ b^{k+1-l} + a^{k+1} + b^{k+1} \\
&= \displaystyle \sum_{l=0}^{k+1} {}_{k+1} C_l \ a^l \ b^{k+1-l}
\end{align*}

よって, $n=k+1$ のときも成り立つ.

(i), \ (ii) より二項定理が示された.


\end{document}
